\documentclass[10pt,a4paper]{amsart}
\usepackage[margin=19mm]{geometry}
\usepackage{amsmath,amssymb,mathtools}
\usepackage[scr=rsfs,cal=euler]{mathalfa}
\usepackage{xfrac}
\usepackage[unicode,hidelinks,bookmarks=false]{hyperref}
\newcommand{\ie}{i.e.\ }
\newcommand{\cf}{cf.\ }
\newcommand{\resp}{respectively\ }
\newcommand{\Subsection}{\subsection}
\providecommand{\nobreakdash}{\nobreak\hbox{-}}
\setlength{\emergencystretch}{2em}
\multlinegap0pt
\usepackage{amssymb}
\usepackage[shortlabels]{enumitem}
\usepackage{cancel}
\newtheorem{lem}{Lemma}
\title{Global existence and time decay for a bipolar Euler-Poisson system with one pressureless and undamped fluid}
\author{Young-Pil Choi, Houzhi Tang, Weiyuan Zou}
\date{Source: arXiv:2609.10959v1 (CC BY 4.0)}
\begin{document}
\maketitle

\noindent\emph{Source and licence.} Global existence and time decay for a bipolar Euler-Poisson system with one pressureless and undamped fluid. Version arXiv:2609.10959v1 (CC BY 4.0), distributed by arXiv under the Creative Commons Attribution 4.0 International licence, http://creativecommons.org/licenses/by/4.0/, as recorded in the arXiv licence metadata for this exact version and shown on the abstract page https://arxiv.org/abs/2609.10959v1. Full licence notice: \texttt{licenses/arxiv-2609.10959-CC-BY-4.0.txt}.\par\smallskip

\noindent\textit{Editorial scope.} Counted unit: the article's lem lem-nonlinear-decay (source0.tex lines 2192--2199). The article's complete original proof of it is delivered with it (source0.tex lines 2201--2408, one \texttt{proof} environment of 208 lines). No mathematics is written, repaired or re-derived: the statement and the proof are the article's own lines, in the article's own notation, and no retained line is substituted or re-flowed.  Retained uncounted as the source-local context the proof needs: lines 283--285; lines 1036--1038; lines 1050--1073; lines 1120--1129; lines 1171--1188; lines 1398--1406; lines 1470--1536; lines 1520--1525; lines 1532--1535; lines 1585--1590; lines 1693--1698; lines 1708--1710; lines 1717--1719; lines 2032--2059; lines 2087--2091.  Nothing was omitted from any retained span, so the package carries no omission record.  No retained line contains a citation, so the wrapper carries no bibliography and no bibliography entry.  No retained line contains a figure environment, an image-inclusion command or drawing code, so no image is delivered and the package declares no asset dependency.  Editorial formatting only, in addition: an AMS \texttt{amsart} preamble that restates the article's own declarations and macros used by the retained lines, namely the environments \texttt{lem} on separate counters, together with the standard packages those lines need.  The retained environments print in this document as Lemma 1 in order of appearance, whereas the article prints the same items with the numbering of its own full document.  The complete native source of the article as distributed in the arXiv e-print for this exact version is bundled unchanged as \texttt{sources/209971/source0.tex}.
\begin{align}\label{main-smallness}
\varepsilon_N:=&\ \|a_{20}\|_{\dot H^{-1}} +\left\|a_{10}-\frac1c a_{20}\right\|_{\dot H^{-1}} +\|a_{10}\|_{H^{N-1}} +\|(u_{10},a_{20},u_{20})\|_{H^N}.
\end{align}

\begin{equation}\label{char-eq}
p(\lambda,\xi) =\det\left(\lambda I+\widehat{\mathcal{L}}(\xi)\right) =\lambda^4+\lambda^3+\left(2+c^2|\xi|^2\right)\lambda^2+\lambda+c^2|\xi|^2.
\end{equation}

\begin{lem}\label{lem-low}
There exists $r_0>0$ sufficiently small such that, for $|\xi|\le 2r_0$, the four roots of \eqref{char-eq} satisfy
\begin{align*}
\lambda_1&=e_0+e_2|\xi|^2+O(|\xi|^4),\quad \lambda_2=j_0+j_2|\xi|^2+O(|\xi|^4),\\
\lambda_3&=z_0+z_2|\xi|^2+O(|\xi|^4),\quad \lambda_4=-c^2|\xi|^2+O(|\xi|^4),
\end{align*}
where $e_0,j_0,z_0$ are the three roots of
\[
\lambda^3+\lambda^2+2\lambda+1=0
\]
and
\begin{align*}
e_2=\frac{-(e_0^2+1)c^2}{4e_0^3+3e_0^2+4e_0+1},\quad
j_2=\frac{-(j_0^2+1)c^2}{4j_0^3+3j_0^2+4j_0+1},\quad
z_2=\frac{-(z_0^2+1)c^2}{4z_0^3+3z_0^2+4z_0+1}.
\end{align*}
Moreover, there exists $\eta_1>0$ such that
\[
\operatorname{Re}\lambda_j(\xi)\le-\eta_1\quad (j=1,2,3),
\qquad
\operatorname{Re}\lambda_4(\xi)\le-\eta_1|\xi|^2,
\quad |\xi|\le2r_0.
\]
\end{lem}

\begin{lem}\label{lem-middle-gap}
	Let $0<r_0<R_0<\infty$ be fixed. There exists a constant
	$ \eta_2 = \eta_2 (c,r_0,R_0)>0$ such that the four roots of
	\eqref{char-eq} satisfy
	\begin{equation}\label{middle-spectral-gap}
		\max_{1\le j\le4}\operatorname{Re}\lambda_j(\xi)
		\le - 2\eta_2, 
		\qquad r_0\le |\xi|\le R_0.
	\end{equation}
\end{lem}

\begin{lem}\label{lem-high}
There exists $R_0>0$ sufficiently large such that, for $|\xi|\ge R_0/2$, the four roots satisfy
\begin{align*}
\lambda_1&=i-\frac{i}{2c^2}|\xi|^{-2}
-\left(\frac12+\frac{i}{8}\right)c^{-4}|\xi|^{-4}+O(|\xi|^{-6}),\\
\lambda_2&=-i+\frac{i}{2c^2}|\xi|^{-2}
-\left(\frac12-\frac{i}{8}\right)c^{-4}|\xi|^{-4}+O(|\xi|^{-6}),\\
\lambda_3&=-\frac12+ic|\xi|+\frac{3i}{8c}|\xi|^{-1}+O(|\xi|^{-2}),\quad \lambda_4 =-\frac12-ic|\xi|-\frac{3i}{8c}|\xi|^{-1}+O(|\xi|^{-2}).
\end{align*}
Hence there exists $\eta_3 >0$ such that
\begin{align}\label{high-eigen-decay}
\begin{aligned}
|e^{\lambda_1t}|+|e^{\lambda_2t}| &\le C\exp\left(-\eta_3  t|\xi|^{-4}\right),\\
|e^{\lambda_3t}|+|e^{\lambda_4t}| &\le Ce^{-\eta_3  t},
\quad |\xi|\ge R_0/2.
\end{aligned}
\end{align}
\end{lem}

\begin{align}\label{low-P4}
P_4=
\begin{pmatrix}
O(|\xi|^2)&O(1)&O(1)&O(1)\\
O(|\xi|^4)&O(|\xi|^2)&O(|\xi|^2)&O(|\xi|^2)\\
O(|\xi|^2)&O(1)&O(1)&O(1)\\
O(|\xi|^4)&O(|\xi|^2)&O(|\xi|^2)&O(|\xi|^2)
\end{pmatrix}.
\end{align}

\begin{lem}[Green-function estimates]\label{lem-green-estimates}
 With the positive constants $\eta_1$, $\eta_2$, and $\eta_3$ from Lemmas \ref{lem-low}, \ref{lem-middle-gap}, and \ref{lem-high}, there exists $C>0$ such that, for all $t\ge0$, the following estimates hold.

\smallskip
\noindent\emph{(i) Low-frequency estimate.} For $|\xi|\le2r_0$,
\begin{align}\label{low-G-pointwise}
|\widehat G_{ij}(t,\xi)|
\le Ce^{-\eta_1  t}+C|\xi|^{\alpha_{ij}}e^{-\eta_1 |\xi|^2t},
\quad 1\le i,j\le4,
\end{align}
and
\[%\begin{align}\label{low-charge-G} 
\left|\left(1,0,-\frac1c,0\right)\widehat G(t,\xi)\right| \le Ce^{-\eta_1  t}(1,1,1,1) +Ce^{-\eta_1 |\xi|^2t} \big(|\xi|^6,|\xi|^4,|\xi|^4,|\xi|^4\big).
\]%\end{align}
If
\[
\widehat{K_\mu^\ell(t)f}(\xi)
=\widehat\chi_\ell(\xi)|\xi|^\mu e^{-\eta_1 |\xi|^2t}\widehat f(\xi),
\quad \mu\ge0,
\]
then, for every integer $k\ge0$,
\begin{align}%\label{low-kernel-L12}
\|\nabla^kK_\mu^\ell(t)f\|_{L^2}
&\le C(1+t)^{-\frac34-\frac{k+\mu}{2}}\|f\|_{L^1}, \notag \\
%\label{low-kernel-L22}
\|\nabla^kK_\mu^\ell(t)f\|_{L^2}
&\le C(1+t)^{-\frac{k+\mu}{2}}\|f\|_{L^2},  \notag \\
\label{low-kernel-L1inf}
\|\nabla^kK_\mu^\ell(t)f\|_{L^\infty}
&\le C(1+t)^{-\frac32-\frac{k+\mu}{2}}\|f\|_{L^1},\\
\label{low-kernel-L2inf}
\|\nabla^kK_\mu^\ell(t)f\|_{L^\infty}
&\le C(1+t)^{-\frac34-\frac{k+\mu}{2}}\|f\|_{L^2}.
\end{align}

\smallskip
\noindent\emph{(ii) Middle-frequency estimate.} For $r_0\le|\xi|\le R_0$,
\begin{align}\label{middle-decay}
|\widehat G_{ij}(t,\xi)|\le Ce^{-\eta_2   t},
\quad 1\le i,j\le4.
\end{align}
Consequently, for every integer $k\ge0$,
\begin{align}\label{middle-G-operator}
\|\nabla^k(G^m(t)*f)\|_{L^2}
+\|\nabla^k(G^m(t)*f)\|_{L^\infty}
\le C_k e^{-\eta_2   t}\|f\|_{L^2}.
\end{align}

\smallskip
\noindent\emph{(iii) High-frequency part.} For $|\xi|\ge R_0/2$,
\begin{align}\label{high-G-pointwise}
|\widehat G_{ij}(t,\xi)|
\le C|\xi|^{-\sigma_{ij}}e^{-\eta_3  t|\xi|^{-4}}
+C|\xi|^{\beta_{ij}}e^{-\eta_3  t},
\quad 1\le i,j\le4.
\end{align}
Furthermore, for every $\ell\ge0$,
\begin{align}\label{regularity-loss-ineq}
e^{-\eta_3  t|\xi|^{-4}}
\le C_\ell(1+t)^{-\frac\ell4}|\xi|^\ell,
\end{align}
and hence, for every integer $k\ge0$ and $\ell\ge1$,
\begin{align}\label{high-linear-general}
\|\nabla^k(G^h(t)*f)\|_{L^2}
\le C_\ell(1+t)^{-\frac\ell4}\|f\|_{H^{k+\ell}}.
\end{align}
\end{lem}

\begin{align}
|\widehat G_{ij}(t,\xi)|
\le C|\xi|^{-\sigma_{ij}}e^{-\eta_3  t|\xi|^{-4}}
+C|\xi|^{\beta_{ij}}e^{-\eta_3  t},
\quad 1\le i,j\le4.
\end{align}

\begin{align}
\|\nabla^k(G^h(t)*f)\|_{L^2}
\le C_\ell(1+t)^{-\frac\ell4}\|f\|_{H^{k+\ell}}.
\end{align}

\begin{align}\label{IBP-Green}
G_{i2}(t)*q_{10}
=\sum_{m=1}^3\partial_mG_{i2}(t)*u_{10,m},\quad
G_{i4}(t)*q_{20}
=\sum_{m=1}^3\partial_mG_{i4}(t)*u_{20,m}.
\end{align}

\begin{align}\label{nonlinear-g}
\begin{aligned}
g_1&=-\operatorname{div}(a_1u_1),\quad g_2=-\mathcal R(a_2)-\operatorname{div}(u_1\cdot\nabla u_1),\quad g_3=-u_2\cdot\nabla a_2-\frac{\gamma-1}{2}a_2q_2,\\ g_4&=\mathcal R(a_2)-\operatorname{div}(u_2\cdot\nabla u_2)
-\frac{\gamma-1}{2}\operatorname{div}(a_2\nabla a_2).
\end{aligned}
\end{align}

\begin{equation}\label{nonlinear-duhamel}
U(t)=G(t)*U_0+\int_0^tG(t-\tau)*F(\tau)\,d\tau.
\end{equation}

\begin{equation}\label{high-bootstrap}
\sup_{0\le\tau\le t}\mathcal E_N(\tau)^{\frac{1}{2}}\le 2\delta
\end{equation}

\begin{lem}\label{lem-nonlinear-source}
Assume \eqref{high-bootstrap}, $\varepsilon_N\le1$, and
$\mathcal M(t)<\infty$. Then, for
$0\le\tau\le t$,
\begin{align}\label{source-L1}
\|F(\tau)\|_{L^1}
&\le C\big(\delta\mathcal M(t)+\mathcal M(t)^2\big)(1+\tau)^{-1}
+C\varepsilon_Ne^{-\frac{\tau}{2}},\\
\label{source-L2}
\|F(\tau)\|_{L^2}
&\le C\big(\delta\mathcal M(t)+\mathcal M(t)^2\big)(1+\tau)^{-\frac{3}{2}},\\
\label{source-H8}
\|F(\tau)\|_{H^8}
&\le C\delta\mathcal M(t)(1+\tau)^{-\frac{5}{4}}.
\end{align}
In addition, with
\[
F_2^{\rm div}:=u_1\cdot\nabla u_1,\quad
F_4^{\rm div}:=u_2\cdot\nabla u_2+\frac{\gamma-1}{2}a_2\nabla a_2,
\]
one has
\begin{align}\label{vector-source-L1L2}
\begin{aligned}
\|F_2^{\rm div}(\tau)\|_{L^1}+\|F_4^{\rm div}(\tau)\|_{L^1} &\le C\big(\delta\mathcal M(t)+\mathcal M(t)^2\big)(1+\tau)^{-1} +C\varepsilon_Ne^{-\frac{\tau}{2}},\\
\|F_2^{\rm div}(\tau)\|_{L^2}+\|F_4^{\rm div}(\tau)\|_{L^2} &\le C\big(\delta\mathcal M(t)+\mathcal M(t)^2\big)(1+\tau)^{-\frac{3}{2}}.
\end{aligned}
\end{align}
\end{lem}

\begin{align}\label{omega2-decay}
\|\omega_2(\tau)\|_{H^2}
\le Ce^{-\frac{\tau}{2}}\|\omega_{20}\|_{H^2}
\le  Ce^{-\frac{\tau}{2}}\|u_{20}\|_{H^3}\le Ce^{-\frac{\tau}{2}}\varepsilon_N.
\end{align}

\begin{lem}\label{lem-nonlinear-decay}
Under \eqref{high-bootstrap}, one has
\begin{align*}%\label{nonlinear-decay-main}
&(1+t)^{\frac{1}{2}}\|(a_1,a_2)(t)\|_{L^2} +(1+t)\|\nabla(a_1,a_2)(t)\|_{H^1} +(1+t)\|(q_1,q_2)(t)\|_{H^2} +(1+t)^{\frac{5}{4}}\mathcal Z(t)\\
&\quad \le C\varepsilon_N+C\delta\mathcal M(t)+C\mathcal M(t)^2,
\end{align*}
where $\varepsilon_N$ is the initial norm in \eqref{main-smallness}.
\end{lem}

\begin{proof}
We estimate the low-, middle-, and high-frequency parts separately.  The constants below are independent of $t$.


\medskip
\noindent\emph{Step 1: low-frequency estimates in $L^2$.} Using the divergence structure of $g_1,g_2$ and $g_4$ in
\eqref{nonlinear-g}, we transfer the outer spatial derivative to the
Green function.
By \eqref{source-L1} and \eqref{vector-source-L1L2},
\begin{align*}
\|F(\tau)\|_{L^1} +\|F_2^{\rm div}(\tau)\|_{L^1} +\|F_4^{\rm div}(\tau)\|_{L^1} \le C\big(\delta\mathcal M(t)+\mathcal M(t)^2\big)(1+\tau)^{-1} +C\varepsilon_Ne^{-\frac{\tau}{2}}.
\end{align*}
Moreover, \eqref{low-G-pointwise} gives the low-frequency orders
\[
\text{rows }1,3:\ (2,0,0,0),
\quad
\text{rows }2,4:\ (4,2,2,2).
\]
Hence Lemma \ref{lem-green-estimates} yields
\begin{align*}
\|(a_1^\ell,a_2^\ell)(t)\|_{L^2} &\le C(1+t)^{-\frac12}\varepsilon_N +C\big(\delta\mathcal M(t)+\mathcal M(t)^2\big) \int_0^t(1+t-\tau)^{-\frac34}(1+\tau)^{-1}\,d\tau\\
&\quad+C\varepsilon_N \int_0^t(1+t-\tau)^{-\frac34}e^{-\frac{\tau}{2}}\,d\tau\\
&\le C(1+t)^{-\frac12} \big(\varepsilon_N+\delta\mathcal M(t)+\mathcal M(t)^2\big),
\end{align*}
where
\[
\int_0^t(1+t-\tau)^{-\frac34}(1+\tau)^{-1}\,d\tau
\le C(1+t)^{-\frac12},
\quad
\int_0^t(1+t-\tau)^{-\frac34}e^{-\frac{\tau}{2}}\,d\tau
\le C(1+t)^{-\frac34}.
\]
For one and two spatial derivatives of the density components,
\begin{align*}
\int_0^t(1+t-\tau)^{-\frac54}(1+\tau)^{-1}\,d\tau  \le C(1+t)^{-1},\quad \int_0^t(1+t-\tau)^{-\frac74}(1+\tau)^{-1}\,d\tau \le C(1+t)^{-1}.
\end{align*}
Therefore,
\[%\begin{align}\label{low-density-nonlinear}
\|\nabla(a_1^\ell,a_2^\ell)(t)\|_{H^1}
\le C(1+t)^{-1}
\big(\varepsilon_N+\delta\mathcal M(t)+\mathcal M(t)^2\big).
\]%\end{align}
For $q_1^\ell,q_2^\ell$, the second and fourth rows contain at least two low-frequency powers, and
\[
\int_0^t(1+t-\tau)^{-\frac74}(1+\tau)^{-1}\,d\tau
\le C(1+t)^{-1}.
\]
Consequently,
\[%\begin{align}\label{low-q-nonlinear}
\|(q_1^\ell,q_2^\ell)(t)\|_{H^2}
\le C(1+t)^{-1}
\big(\varepsilon_N+\delta\mathcal M(t)+\mathcal M(t)^2\big).
\]%\end{align}


\medskip
\noindent\emph{Step 2: middle-frequency estimate.}
On the fixed annulus $r_0\le|\xi|\le R_0$, Lemma \ref{lem-green-estimates} gives an exponential semigroup bound.  Since all derivatives are equivalent on this annulus,
\begin{align}\label{middle-nonlinear}
\begin{aligned}
\|U^m(t)\|_{H^3} &\le Ce^{-\eta_2   t}\|U_0\|_{L^2} +C\int_0^te^{-\eta_2  (t-\tau)}\|F(\tau)\|_{L^2}\,d\tau \\
&\le C(1+t)^{-\frac{3}{2}} \big(\varepsilon_N+\delta\mathcal M(t)+\mathcal M(t)^2\big),
\end{aligned}
\end{align}
where \eqref{source-L2} was used in the last line.

\medskip
\noindent\emph{Step 3: high-frequency estimate.} Choose $\ell=5$ in \eqref{regularity-loss-ineq}.  From \eqref{high-G-pointwise}, for $0\le k\le3$,
\[ 
\left\|\nabla^k\int_0^tG^h(t-\tau)*F(\tau)\,d\tau\right\|_{L^2}
\le C\int_0^t(1+t-\tau)^{-\frac{5}{4}}\|F(\tau)\|_{H^{k+5}}\,d\tau.
\] 
Since $k+5\le8$, Lemma \ref{lem-nonlinear-source} yields
\[
\left\|\int_0^tG^h(t-\tau)*F(\tau)\,d\tau\right\|_{H^3} \le C\delta\mathcal M(t) \int_0^t(1+t-\tau)^{-\frac{5}{4}}(1+\tau)^{-\frac{5}{4}}\,d\tau \le C\delta\mathcal M(t)(1+t)^{-\frac{5}{4}}.
\]
The homogeneous high-frequency term is estimated by \eqref{high-linear-general}.  Since {$N\geq 10$} and $q_{i0}=\operatorname{div}u_{i0} (i=1,2)$,
\[
\|U_0\|_{H^8}\le C\varepsilon_N.
\]
Therefore
\begin{align}\label{high-nonlinear-final}
\|U^h(t)\|_{H^3}
\le C(1+t)^{-\frac{5}{4}}\big(\varepsilon_N+\delta\mathcal M(t)\big).
\end{align}


Combining Steps 1--3, we obtain
\begin{align}
\|(a_1,a_2)(t)\|_{L^2}
&\le C(1+t)^{-\frac12}\big(\varepsilon_N+\delta\mathcal M(t)+\mathcal M(t)^2\big),\label{T01}\\
\|\nabla(a_1,a_2)(t)\|_{H^1}
+\|(q_1,q_2)(t)\|_{H^2}
&\le C(1+t)^{-1}\big(\varepsilon_N+\delta\mathcal M(t)+\mathcal M(t)^2\big).\label{T02}
\end{align}


\medskip
\noindent\emph{Step 4: the $W^{1,\infty}$ norm.} We first estimate the low-frequency density components. By \eqref{low-P4}, the slow mode in the first and third rows contains two powers of $|\xi|$ on the initial datum $a_{10}$, but contains no additional low-frequency factor on $a_{20}$. On the other hand, $q_{10}=\operatorname{div}u_{10}$ and $q_{20}=\operatorname{div}u_{20}$ each provide one power of $|\xi|$ through \eqref{IBP-Green}. It follows from the low-frequency $L^2$--$L^\infty$ multiplier estimate that
\[
\left\| \bigl(G^\ell(t)*U_0\bigr)_{(1,3)} \right\|_{W^{1,\infty}} \le C\int_{|\xi|\le 2r_0} e^{-\eta_1 |\xi|^2t} |\widehat a_{20}(\xi)|\,d\xi +C\varepsilon_N(1+t)^{-\frac54}.
\]
Here $(\cdot)_{(1,3)}$ denotes the first and third components. Indeed, the contribution of $a_{10}$ decays at least as $(1+t)^{-7/4}$, while the contributions of $q_{10}$ and $q_{20}$ decay at least as $(1+t)^{-5/4}$. The exponentially stable low-frequency modes are also bounded by $C\varepsilon_N(1+t)^{-5/4}$.

We next estimate the Duhamel term. Applying the full Duhamel formula \eqref{nonlinear-duhamel}, and using \eqref{IBP-Green}, Bernstein's inequality, \eqref{low-kernel-L1inf}--\eqref{low-kernel-L2inf}, and \eqref{source-L1}--\eqref{source-L2}, we obtain
\begin{align*}
\|(a_1^\ell,a_2^\ell)(t)\|_{W^{1,\infty}} &\le C\int_{|\xi|\le 2r_0} e^{-\eta_1 |\xi|^2t} |\widehat a_{20}(\xi)|\,d\xi +C\varepsilon_N(1+t)^{-\frac54} 
+C\varepsilon_N \int_0^t (1+t-\tau)^{-\frac32}e^{-\frac{\tau}{2}}\,d\tau\\
&\quad +C\bigl(\delta\mathcal M(t)+\mathcal M(t)^2\bigr) \left[ \int_0^{\frac{t}{2}} (1+t-\tau)^{-\frac32}(1+\tau)^{-1}\,d\tau +\int_{\frac{t}{2}}^t (1+t-\tau)^{-\frac34}(1+\tau)^{-\frac32}\,d\tau \right].
\end{align*}
The negative Sobolev assumption on $a_{20}$ gives
\[
 \int_{|\xi|\le 2r_0} e^{-\eta_1 |\xi|^2t} |\widehat a_{20}(\xi)|\,d\xi  \le \left( \int_{|\xi|\le 2r_0} |\xi|^2e^{-2\eta_1 |\xi|^2t}\,d\xi \right)^{\frac12} \|a_{20}\|_{\dot H^{-1}} \le C(1+t)^{-\frac54} \|a_{20}\|_{\dot H^{-1}}.
\]
Moreover,
\begin{align*}
\int_0^t (1+t-\tau)^{-\frac32}e^{-\frac{\tau}{2}}\,d\tau &\le C(1+t)^{-\frac32},\\
\int_0^{\frac{t}{2}} (1+t-\tau)^{-\frac32}(1+\tau)^{-1}\,d\tau &\le C(1+t)^{-\frac54},\\
\int_{\frac{t}{2}}^t (1+t-\tau)^{-\frac34}(1+\tau)^{-\frac32}\,d\tau &\le C(1+t)^{-\frac54}.
\end{align*}
Combining these estimates yields
\[
\|(a_1^\ell,a_2^\ell)(t)\|_{W^{1,\infty}} \le C(1+t)^{-\frac54} \bigl( \varepsilon_N +\delta\mathcal M(t) +\mathcal M(t)^2 \bigr).
\]
 
Since $\operatorname{curl}u_1=0$, we get $u_1^\ell=-\nabla(-\Delta)^{-1}q_1^\ell$. Applying $\nabla(-\Delta)^{-1}$ directly to the second row of the full
Duhamel formula \eqref{nonlinear-duhamel}, the low-frequency orders
$(4,2,2,2)$ in \eqref{low-G-pointwise} become $(3,1,1,1)$.  Hence,
by \eqref{low-kernel-L1inf}--\eqref{low-kernel-L2inf},
\eqref{source-L1}--\eqref{source-L2}, and a splitting at $t/2$,
\begin{align*}
\|u_1^\ell(t)\|_{W^{1,\infty}}
&\le C(1+t)^{-\frac54}\varepsilon_N
+C\int_0^{\frac{t}{2}}(1+t-\tau)^{-2}\|F(\tau)\|_{L^1}\,d\tau\\
&\quad+C\int_{\frac{t}{2}}^t(1+t-\tau)^{-\frac54}\|F(\tau)\|_{L^2}\,d\tau
+C\varepsilon_Ne^{-\eta_1  t}
+C\int_0^te^{-\eta_1 (t-\tau)}\|F(\tau)\|_{L^2}\,d\tau\\
&\le C(1+t)^{-\frac54}
\big(\varepsilon_N+\delta\mathcal M(t)+\mathcal M(t)^2\big),
\end{align*}
where the exponentially stable part satisfies the same bound since
$\widehat\chi_\ell(\xi)|\xi|^{-1}\in L^2_\xi$.  Thus
\[%\begin{align}\label{low-u1-W1inf}
\|u_1^\ell(t)\|_{W^{1,\infty}}
\le C(1+t)^{-\frac54}
\big(\varepsilon_N+\delta\mathcal M(t)+\mathcal M(t)^2\big).
\]%\end{align}
 

For the second velocity, set $\omega_2=\operatorname{curl}u_2$. By the Helmholtz decomposition,
\[
u_2=-\nabla(-\Delta)^{-1}q_2
+\operatorname{curl}(-\Delta)^{-1}\omega_2.
\]
Since the low-frequency cut-off commutes with these Fourier multipliers,
taking the low-frequency part gives
\[
u_2^\ell=-\nabla(-\Delta)^{-1}q_2^\ell
+\operatorname{curl}(-\Delta)^{-1}\omega_2^\ell.
\]
For the first term, the fourth row of the low-frequency Green matrix
contains at least two powers of $|\xi|$, while
$\nabla(-\Delta)^{-1}$ contributes one factor $|\xi|^{-1}$.
Hence the same low-frequency kernel estimates as above yield
\[
\|\nabla(-\Delta)^{-1}q_2^\ell(t)\|_{W^{1,\infty}}
\le C(1+t)^{-\frac54}
\big(\varepsilon_N+\delta\mathcal M(t)+\mathcal M(t)^2\big).
\]
For the curl part, \eqref{omega2-decay} gives
\begin{align*}
\|\operatorname{curl}(-\Delta)^{-1}\omega_2^\ell(t)\|_{W^{1,\infty}}
&\le C\|\omega_2(t)\|_{H^2}
\le C\varepsilon_Ne^{-\frac{t}{2}}.
\end{align*}
Therefore,
\[%\begin{align}\label{low-u2-W1inf}
\|u_2^\ell(t)\|_{W^{1,\infty}}
\le C(1+t)^{-\frac54}
\big(\varepsilon_N+\delta\mathcal M(t)+\mathcal M(t)^2\big).
\]%\end{align}

On middle frequencies, the div--curl decomposition, \eqref{middle-nonlinear}, and
\eqref{omega2-decay} yield
\begin{align*}
 \|(a_1^m,a_2^m,u_1^m,u_2^m)(t)\|_{W^{1,\infty}}  \le C\|U^m(t)\|_{H^3}+C\|\omega_2^m(t)\|_{H^2} \le C(1+t)^{-\frac54}
\big(\varepsilon_N+\delta\mathcal M(t)+\mathcal M(t)^2\big).
\end{align*}
For the high-frequency part,
\begin{align*}
\|\nabla u_1^h\|_{H^2}\le C\|q_1^h\|_{H^2},\quad
\|\nabla u_2^h\|_{H^2}\le C\big(\|q_2^h\|_{H^2}+\|\omega_2^h\|_{H^2}\big),
\end{align*}
so \eqref{high-nonlinear-final}, \eqref{omega2-decay}, and
$H^3\hookrightarrow W^{1,\infty}$ imply
\[%\begin{align}\label{high-W1inf}
\|(a_1^h,a_2^h,u_1^h,u_2^h)(t)\|_{W^{1,\infty}}
\le C(1+t)^{-\frac54}
\big(\varepsilon_N+\delta\mathcal M(t)+\mathcal M(t)^2\big).
\]%\end{align}
Combining the three frequency regions gives
\begin{align}\label{Z-decay}
\mathcal Z(t)
\le C(1+t)^{-\frac54}
\big(\varepsilon_N+\delta\mathcal M(t)+\mathcal M(t)^2\big).
\end{align}
Combining \eqref{T01}, \eqref{T02}, and \eqref{Z-decay} completes the proof.
\end{proof}

\end{document}
